\section{Deployment and Operational Support}
\label{sec:deployment-support}

\subsection{Operational and Environmental Constraints}\label{subsec:operational-and-environmental-constraints}
For subterranean irrigation deployments, the physical environment and service model materially change requirements.
We therefore set explicit operational constraints that must be observed when deploying TrackOne artifacts in the field:
\begin{itemize}[nosep]
  \item \textbf{Deployment tolerance:} Pods are expected to be wet, buried, and inaccessible for long durations; hardware must therefore favour robustness and conservative scheduling for maintenance windows.
  \item \textbf{Maximum allowed downtime:} Configurable per site (typical value 24 hours); longer outages must be escalated to human operators.
  \item \textbf{Acceptable data loss:} Define a site-specific bound (e.g., no more than 0.5\% of day's facts lost) beyond which operator rollback or manual audit is required.
  \item \textbf{Sampling frequency:} Default 1--10 samples per hour per pod depending on sensor; higher rates must be justified by an explicit risk/benefit analysis.
  \item \textbf{Firmware update policy:} Updates must support atomic rollback and require an operator-confirmed maintenance window; over-the-air auto-update that can change actuation logic is out-of-scope for TrackOne.
  \item \textbf{Storage limits:} On-device and gateway storage quotas must be set; when storage is full, pods/gateway must stop non-essential writes and raise an alert to operator channels rather than silently dropping telemetry.
  \item \textbf{Stationary calendar usage:} Use of the stationary calendar in \ac{CI} and weekly ``ratchet'' jobs is MANDATORY when public \ac{OTS} pools are unstable.
  \item \textbf{Embedded viability:} At least one Cortex-M0+/M4 and one RISC-V firmware build \ac{SHA}LL be kept
  in-tree and exercised regularly (stress harness or FPGA target) to guarantee that TrackOne-core behaves identically
  across \ac{MCU} families.
  \item \textbf{GIS-constrained placement:} Sensor pod placements \ac{SHA}LL be drawn from pre-approved \ac{GIS} polygons (see Section~\ref{sec:site_planning}), not arbitrary coordinates.
\end{itemize}

\subsection{GIS-Informed Pod Placement Constraints}\label{subsec:gis-informed-pod-placement-constraints}
For subterranean irrigation corridors (khettara-like tunnels) and adjacent heritage structures, pod placement must be planned with basic geospatial constraints that can be encoded in deployment runbooks:
In water-domain deployments, the site plan and placement constraints also support governance and reporting obligations (Section~\ref{subsec:water-law}).
\begin{itemize}[nosep]
  \item \textbf{Buffer zones:} Maintain minimum stand-off distances from protected masonry, inscriptions, and water intakes (site-specific, e.g., 30--50~cm) to preserve reversibility and avoid damage.
  \item \textbf{Radio propagation:} Prefer line-of-sight along shafts; where bends or thick walls exist, introduce relay pods at bends mapped in \ac{GIS} to keep link budgets within margins.
  \item \textbf{Hydrology overlays:} Avoid low points identified by flood maps for primary pods; place flood sentinels explicitly at risk points to trigger maintenance/evacuation procedures.
  \item \textbf{Access waypoints:} Every pod should be within a serviceable reach from registered maintenance paths; the \ac{GIS} plan must list approach routes and safe access times (seasonal constraints).
  \item \textbf{Cultural heritage constraints:} Mounts must be non-invasive (adhesive/straps) and removable; no drilling or permanent fixtures without curator approval.
\end{itemize}

\subsection{Spatial Database Requirements}\label{subsec:spatial-database-requirements}
To keep \ac{GIS} reasoning close to the verifiable pipeline while avoiding heavyweight external services, TrackOne MAY use a
Spatialite-backed SQLite database for pod placement, maintenance routes, and hydrology overlays.
We require:

\begin{description}[style=nextline]
  \item[SD-1 Single source of spatial truth] Site polygons, buffer zones, and approved maintenance routes \ac{SHA}LL be
  stored in a version-controlled Spatialite database.
  Generated pod placement plans MUST reference this database
  rather than ad-hoc shapefiles.
  \item[SD-2 Read-side separation] The ingest and ledger paths \ac{SHA}LL treat the Spatialite database as a read-only state.
  Failure or corruption of spatial indices MUST NOT compromise ledger integrity; it can only degrade planning and
  visualization.
  \item[SD-3 Export compatibility] Spatialite contents \ac{SHA}LL be exportable to standard \ac{GIS} formats (e.g., GeoPackage,
  Shapefile) for use in Q\ac{GIS}/Arc\ac{GIS} without loss of topology or metadata relevant to pod placement.
\end{description}

\subsection{Resilience Requirements}\label{subsec:resilience-requirements}
We add concise, testable targets that express how the pipeline maintains useful service across outages and partial failures:
\begin{description}[style=nextline]
  \item[RR-1 Anchor backlog clearance] After an uplink outage of up to 7 days, gateways shall clear the \ac{OTS} anchoring backlog within 48~hours of connectivity restoration while preserving per-day artifact immutability.
  \item[RR-2 Multi-gateway replication] In multi-gateway deployments, day artifacts shall be replicated to at least one peer within 24~hours (\ac{ADR}-019 Supernode Train) to bound single-node loss.
  \item[RR-3 Stationary calendar fallbacks] When public calendars are unavailable, stationary calendar proofs shall be accepted with a visible \emph{pending public anchor} status in meta until upgraded.
  \item[RR-4 Peer attestation quorum] Day roots shall collect at least $q$ operator signatures (default $q{=}2$ of $n$) before being marked \emph{attested} in operator dashboards.
  \item[RR-5 GIS-linked maintenance] Pods flagged by flood sentinel events or repeated RF failures must enter a \ac{GIS} task list within 24~hours, including route and safety notes for field teams.
\end{description}

\subsection{Energy Budget and Coin-Cell Viability}\label{subsec:energy-budget-and-coin-cell-viability}
For buried pods powered by small primary cells (e.g., CR2032) or harvesters plus a storage element, we explicitly
separate average energy use (mAh / Wh) from peak current viability (ability to sustain \ac{LoRa} TX bursts):
\begin{itemize}[nosep]
  \item \textbf{Assumed cell:} CR2032-class, $\approx 200$~mAh usable at low drain, $E \approx 0.6$~Wh at 3.0~V.
  \item \textbf{Mission duration:} 6~months ($\approx 4320$~h).
  \item \textbf{Sleep current target:} $I_{\text{sleep}} \le 5\,\mu\text{A}$ (\ac{MCU} + leakage + powered-down sensors).
  \item \textbf{Uplink cadence:} 1--4 uplinks/day, each carrying a compact summary (typically $< 80$~B).
  \item \textbf{LoRa TX envelope:} $I_{\text{tx}} \in [40,120]$~mA, airtime 0.5--2.0~s per uplink depending on \ac{SF}/payload.
\end{itemize}
We use the following duty-cycle model for capacity estimates:
\begin{equation}
  E_{\text{wake}} \approx V \cdot (I_{\text{mcu}} t_{\text{mcu}} + I_{\text{sensor}} t_{\text{sensor}} + I_{\text{tx}} t_{\text{tx}} + I_{\text{rx}} t_{\text{rx}}),\label{eq:equation2}
\end{equation}
with total charge over the mission
\begin{equation}
  Q_{\text{total}} \approx I_{\text{sleep}} T_{\text{mission}} + N_{\text{wake}} \cdot \frac{E_{\text{wake}}}{V}.\label{eq:equation}
\end{equation}
Under representative parameters (sleep 5~$\mu$A, 1--4 uplinks/day, $t_{\text{tx}} \le 2$~s) we obtain $Q_{\text{total}} \approx
30$--65~mAh over 6~months: \emph{feasible on a single CR2032 by energy}.
However, peak current and voltage sag during
\ac{LoRa} TX require a reservoir capacitor (or supercap) and in practice motivate either two cells in parallel or a
harvester + storage element.

We therefore state explicit, testable requirements:
\begin{description}[style=nextline]
  \item[ER-1 Coin-cell budget] A reference pod configuration (\ac{MCU} + SHTC3-only sensing, 1 uplink/day, $t_{\text{tx}} \le 1$~s at
  moderate \ac{SF}) \ac{SHA}LL demonstrate a measured average current $I_{\text{avg}} \le 20\,\mu$A over a 7-day endurance test.
  \item[ER-2 Peak-current mitigation] Pods using CR2032-class cells MUST include a local reservoir capacitor sized to
  keep supply voltage within the \ac{MCU} and radio operating bounds during \ac{LoRa} TX peaks.
  Designs that cannot meet this
  \ac{SHA}LL use a different cell chemistry or an energy-harvesting front-end (e.g., BQ25570 + supercap).
  \item[ER-3 Reporting of assumptions] Deployment runbooks \ac{SHA}LL record assumed uplink cadence, \ac{SF},
  and worst-case airtime so that operators can relate conditions back to the energy model.
\end{description}

\subsection{Traceability and Verification}
\label{subsec:requirements-traceability}

Requirements are verified via specific automated tests in the \ac{CI} pipeline:

\begin{table}[H]
\centering
\small
\caption{Requirements Traceability Matrix (Functional)}
\label{tab:req-traceability-fun}
\begin{tabularx}{\textwidth}{|l|>{\raggedright\arraybackslash}X|l|}
\hline
\textbf{Req ID} & \textbf{Requirement Description} & \textbf{Verification Artifact / Test} \\ \hline
FR-1  & Ingest sensor frames           & \texttt{test\_pod\_sim}, \texttt{test\_framed\_ingest} \\ \hline
FR-2  & Anti-replay window             & \texttt{test\_replay\_window\_unit} \\ \hline
FR-3  & Ledger writing                 & \texttt{test\_merkle\_batcher\_sees\_facts} \\ \hline
FR-4  & Audit logging                  & \texttt{test\_framed\_ingest} (log inspection) \\ \hline
FR-5  & Merkle batching                & \texttt{test\_merkle\_determinism} \\ \hline
FR-6  & \ac{OTS} stamping                   & \texttt{test\_ots\_integration} \\ \hline
FR-7  & Verification tooling           & \texttt{verify\_cli} (End-to-End) \\ \hline
FR-8  & Config management              & \texttt{test\_gateway\_pipeline} \\ \hline
FR-9  & Error handling                 & \texttt{test\_framed\_security} \\ \hline
FR-10 & Stationary anchor              & \texttt{tests/e2e/test\_ots\_calendar\_pipeline} \\ \hline
FR-11 & Replay invariants              & \texttt{tests/integration/test\_replay\_merkle} \\ \hline
FR-12 & Cross-impl equivalence         & \texttt{tests/integration/test\_rust\_core\_equiv} \\ \hline
FR-13 & Multi-\ac{ISA} determinism          & Rust suite (ARM/RISC-V) \\ \hline
FR-14 & Deterministic commitments       & \texttt{crates/trackone-core/src/cbor.rs} \\ \hline
FR-15 & Deterministic time             & \texttt{tests/integration/test\_sensorthings} \\ \hline
\end{tabularx}
\end{table}

\begin{table}[H]
\centering
\small
\caption{Requirements Traceability Matrix (Non-Functional and Constraints)}
\label{tab:req-traceability-non-fun}
\begin{tabularx}{\textwidth}{|l|>{\raggedright\arraybackslash}X|l|}
\hline
\textbf{Req ID} & \textbf{Requirement Description} & \textbf{Verification Artifact / Test} \\ \hline
NFR-1 & Integrity                      & \texttt{verify\_cli} (End-to-End) \\ \hline
NFR-2 & Concerns split                 & \texttt{test\_merkle\_batcher\_sees\_duplicate\_free} \\ \hline
NFR-3 & Performance                    & Benchmarks in the cryptographic design chapter \\ \hline
NFR-4 & Dependability                  & \ac{CI} Stress tests \\ \hline
NFR-5 & Testability                    & This matrix and 85\% coverage \\ \hline
NFR-6 & Maintainability                & \ac{ADR} inventory and appendix support material \\ \hline
NFR-7 & Anchoring resilience           & \ac{CI} ratchet (stationary calendar) \\ \hline
NFR-8 & Implementation consistency      & Rust + Python golden corpus \\ \hline
C-1   & Platform (Linux/Python)        & \ac{CI} Environment \\ \hline
C-2   & Cryptography (Standard)        & \texttt{test\_crypto\_vectors} \\ \hline
C-5   & Logging (No secrets)           & Bandit scan (\ac{ADR}-009) \\ \hline
ER-1  & Coin-cell budget               & 7-day endurance test report \\ \hline
ER-2  & Peak-current mitigation        & Hardware review / traces \\ \hline
R-1   & Anchoring delay                & \texttt{test\_ots\_retry\_logic} \\ \hline
R-2   & Data loss limit                & Operational monitoring alerts \\ \hline
R-5   & Auditability                   & Merkle-anchored action metadata tests \\ \hline
SA-1/2& SensorThings                   & \texttt{tests/integration/test\_sensorthings} \\ \hline
SD-1/2& Spatialite                     & \texttt{tests/integration/test\_spatialite\_site\_db} \\ \hline
\end{tabularx}
\end{table}

\section{Governance and Publication Notes}
\label{sec:governance-notes}

\subsection{Documentation and traceability}\label{subsec:documentation-and-traceability}
For an end-of-year project intended to be deployed in challenging environments, documentation is not optional.
The following artifacts are ``living'' and must be kept in sync with code and deployments:
\begin{itemize}[nosep]
  \item Requirements (including the irrigation-specific requirements added in this paper) with traced links to tests and \ac{ADR}s.
  \item Architecture diagrams and deployment profiles (per-site) that document pod/gateway/backend interactions and escalation paths.
  \item Test plans and runbooks covering on-pod diagnostics, failure-injection scenarios (link loss, clock drift, storage exhaustion), and recovery procedures.
  \item Known limitations and mitigation checklists attached to each release (what an operator must not do with this release).
\end{itemize}

These artifacts should be versioned alongside software releases and available to field operators and auditors.

\subsection{Operational procedures}\label{subsec:operational-procedures}
Deploying TrackOne in irrigation contexts requires clear operational playbooks.
Minimal recommended procedures:
\begin{enumerate}[nosep]
  \item Site bootstrap: verify provisioning (device table v1.0), perform a 72-hour pilot with local monitoring enabled, and confirm anchor delivery for two consecutive days before enabling any automatic actuation that affects irrigation.
  \item Firmware update: stage updates to a pilot subset of devices for 7 days; monitor telemetry integrity and rollback on anomalous behavior.
  \item Storage exhaustion: when on-device or gateway storage reaches 80\% capacity, system must stop non-critical writes and alert operators; at 95\% capacity it must trigger emergency archival and operator intervention.
  \item Operator override logging: any manual override must create an auditable record that is included in the next day's Merkle root and proof, so that later audits can reconstruct operator interventions.
\end{enumerate}

\subsection{Ethical and regulatory considerations}\label{subsec:ethical-and-regulatory-considerations}

\textbf{Fact confidentiality:} Telemetry payloads remain local (gateway storage).\ Only Merkle roots are public (Bitcoin blockchain via \ac{OTS}). Individual facts cannot be reconstructed from blockchain data.

\textbf{Right to be forgotten (GDPR):} Local fact deletion is possible (gateway operator decision).\ However, Merkle root in Bitcoin blockchain is immutable.\ Heritage sites should assess \ac{GDPR} applicability before deployment.

\textbf{Metadata leakage:} Frame timing and pod IDs are visible to \ac{LoRa} eavesdroppers.\ XChaCha20 encryption hides payload content but not traffic patterns.

\textbf{Open standards and reuse:} Publishing a SensorThings-compatible \ac{API} means heritage authorities and researchers can consume TrackOne telemetry with off-the-shelf clients; \ac{GIS} and hydrology feeds can be joined without bespoke adapters.
The canonical TrackOne facts and \ac{OTS}-anchored Merkle roots underpin that \ac{API}, so the public view is a projection over an immutable ledger rather than a mutable database without provenance.

\subsection{Environmental Impact}\label{subsec:environmental-impact}

\textbf{Battery disposal:} CR2032 lithium cells require proper recycling.\ Multi-year autonomy reduces disposal frequency vs.\ annually-replaced batteries.

\textbf{Bitcoin energy consumption:} \ac{OTS} leverages existing Bitcoin infrastructure (no additional mining).\ Daily batching minimizes per-fact environmental impact (1 anchor amortized across 1000+ facts).

\subsection{Heritage Site Stewardship}\label{subsec:heritage-site-stewardship}

\textbf{Reversibility:} Adhesive pod mounts are non-invasive and removable without surface damage.\ Critical for heritage site preservation ethics.

\textbf{Curator control:} Gateway operator decides what data to collect and publish (Merkle roots).\ Public verification enables accountability but respects site autonomy.

\textbf{Long-term access:} Open-source design ensures data remains accessible even if TrackOne project discontinues.\ Bitcoin verification survives organizational changes.

\subsection{Responsible Disclosure}\label{subsec:responsible-disclosure}

\textbf{Security vulnerabilities:} All cryptographic designs documented in \ac{ADR}s. Independent security review encouraged (MIT license).\ No security-by-obscurity.

\textbf{Test coverage:} ~85\% coverage with 181 tests provides confidence but not proof of correctness.

\textbf{Known attack vectors:} Physical pod tampering, gateway compromise, replay attacks beyond window size.\ Documented in \ac{ADR}-002 threat model.

\subsection{Stakeholder expectations and ethical scope management}\label{subsec:stakeholder-expectations-and-ethical-scope-management}
Field stakeholders occasionally push comfort-zone technologies (e.g., generic image watermarking or steganography) even when they do not match the telemetry problem space.
Heritage pods emit scalar facts; there are no images today.
Implementing watermarking to appease expectations would misrepresent delivered functionality and divert scarce energy budget.
Instead, the project documents alternative tracks (Option B QIM watermarking, Option C PyO3/Rust port) in \ac{ADR}s, clearly labeling them as future work contingent on hardware funding and bench instrumentation.

Equipment discussions are treated transparently: new hardware (oscilloscope, FPGA stimulus%, bio-impedance rigs
) must map to a concrete experiment plan tied to documented milestones.
Until funding is committed, TrackOne prioritizes the achievable Option A pipeline and avoids promising features (``watermark'' keywords) that are not implemented.
This preserves integrity of reporting and maintains trust with community partners who rely on verifiable telemetry rather than speculative features.

\subsection{Internet-Draft Submission}
\label{subsec:internet-draft-submission}

The TrackOne ledger profile is also documented as an \ac{IETF} Internet-Draft to make the artifact format, verification semantics, and deployment assumptions reviewable outside this report.

\subsubsection{Scope and Positioning}
The draft is informational and focuses on the verifiable telemetry ledger profile: canonical facts, deterministic commitment bytes, Merkle batching, proof sidecars, and verifier expectations.
It does not attempt to standardize deployment-specific hardware, curator workflows, or site planning.
Within the report structure, the draft should be read as publication strategy and protocol-positioning support for Option A rather than as a separate implementation track.

\subsubsection{Implementation Status}
The draft reflects the implemented TrackOne pipeline: Python gateway logic, Rust TrackOne-core parity work, deterministic test corpora, OpenTimestamps integration, and read-only projection support.
Its value is to expose the contract for external review while keeping the main report focused on system design and evaluation.

\subsubsection{Disclosure}
The draft should continue to disclose implementation maturity clearly:
\begin{itemize}[nosep]
  \item implemented and tested behaviors that exist in the repository;
  \item hardening items still pending within Track One;
  \item later-track concepts that are deliberately out of scope for the current release.
\end{itemize}
